\chapter{Digital Signature: Digital Signature Algorithm (DSA)}
\label{ch:dsa}

\section{Historical Background, Development, and Design Motivation}
In the late 1980s, the widespread adoption of RSA for digital signatures raised critical legal and governmental challenges. RSA was encumbered by commercial patents held by RSA Data Security Inc. Furthermore, the United States government sought a public-key standard dedicated strictly to digital authentication and non-repudiation, intentionally precluding general message encryption to comply with Cold War-era cryptographic export control regulations (such as ITAR).

In August 1991, the National Institute of Standards and Technology (NIST) proposed the \textbf{Digital Signature Algorithm (DSA)} as part of the Digital Signature Standard (DSS), formally standardized in 1994 as FIPS PUB 186 and updated in FIPS PUB 186-4 \cite{fips186_4}. DSA was designed by David Kravitz (a former NSA cryptographer), adapting the mathematical structures of the ElGamal and Schnorr signature schemes over finite-field cyclic subgroups. By operating within a prime-order subgroup of $\mathbb{Z}_p^*$, DSA produced signatures that were significantly more compact than RSA signatures (yielding two 160-bit integers, totaling 320 bits, compared to a 1024-bit RSA signature).

\section{Underlying Mathematical Foundations and Operating Principles}
DSA operations take place within a prime-order subgroup of the multiplicative group of integers modulo a prime $p$, denoted $\mathbb{Z}_p^*$.

\subsection{Cyclic Subgroup Structure and Generator Selection}
Let $p$ and $q$ be primes such that $q$ divides $(p - 1)$ (denoted $q \mid (p - 1)$). By Lagrange's Theorem, the multiplicative group $\mathbb{Z}_p^*$ (which has order $p - 1$) contains a unique cyclic subgroup $G_q$ of prime order $q$.
To construct a generator $g$ of this subgroup $G_q$:
\begin{enumerate}
    \item Select an arbitrary integer $h \in \{2, \dots, p - 2\}$.
    \item Compute:
    \begin{equation}
    g = h^{(p - 1)/q} \bmod p
    \end{equation}
    \item If $g > 1$, then $g$ generates $G_q$, satisfying:
    \begin{equation}
    g^q \equiv (h^{(p - 1)/q})^q \equiv h^{p - 1} \equiv 1 \pmod p \quad (\text{by Fermat's Little Theorem})
    \end{equation}
\end{enumerate}

\subsection{The Discrete Logarithm Problem (DLP)}
Given the subgroup generator $g$ and an element $y \in G_q$, finding an integer $x \in \mathbb{Z}_q$ such that:
\begin{equation}
y \equiv g^x \pmod p
\end{equation}
is the Discrete Logarithm Problem (DLP). While modular exponentiation $g^x \bmod p$ can be computed efficiently in $\mathcal{O}(\log x)$ steps via square-and-multiply, the inverse problem has no known polynomial-time classical solution.

\subsection{Mathematical Proof of Verification Correctness}
The signature verification checks whether $v \equiv r \pmod q$.
During signing:
\begin{equation}
s \equiv k^{-1} (H(M) + x \cdot r) \pmod q
\end{equation}
Multiplying both sides by $k \cdot s^{-1} \equiv k \cdot w \pmod q$:
\begin{equation}
k \equiv s^{-1} H(M) + s^{-1} x \cdot r \equiv w \cdot H(M) + w \cdot r \cdot x \equiv u_1 + u_2 \cdot x \pmod q
\end{equation}
Since $g$ has order $q$ in $\mathbb{Z}_p^*$, any exponent reduction modulo $q$ is valid:
\begin{equation}
v = (g^{u_1} y^{u_2} \bmod p) \bmod q = (g^{u_1} (g^x)^{u_2} \bmod p) \bmod q = (g^{u_1 + u_2 x} \bmod p) \bmod q
\end{equation}
Substituting $u_1 + u_2 x \equiv k \pmod q$:
\begin{equation}
v \equiv (g^k \bmod p) \bmod q \equiv r \pmod q
\end{equation}
This confirms that an untampered signature produced with private key $x$ will always satisfy $v = r$.

\section{Complete Operational Procedure}
The complete operational lifecycle of DSA comprises Domain Parameter and Keypair Generation, Signature Generation by the sender, and Signature Verification by the recipient, as illustrated in Figure \ref{fig:dsa_flowchart}.

\begin{figure}[htbp]
\centering
\begin{tikzpicture}[>=stealth,
    block/.style={rectangle, draw=blue!70!black, fill=blue!5, thick, text width=8.5cm, align=center, rounded corners=3pt, inner sep=7pt},
    line/.style={draw, -latex, thick}]
    
    \node (sign) [block] {\textbf{Digital Signature Generation Phase:}\\
        \vspace{2pt}
        \small Ephemeral nonce $k \in_R [1, q-1]$\\
        $r = (g^k \bmod p) \bmod q$\\
        $s = k^{-1}(H(M) + x \cdot r) \bmod q$\\
        $\implies \textbf{Signature Output: } (r, s)$};
    \node (ver) [block, below=0.7cm of sign] {\textbf{Digital Signature Verification Phase:}\\
        \vspace{2pt}
        \small Verify $0 < r, s < q$, then compute $w = s^{-1} \bmod q$\\
        $u_1 = H(M)w \bmod q, \quad u_2 = rw \bmod q$\\
        $v = ((g^{u_1} y^{u_2}) \bmod p) \bmod q$\\
        $\implies \textbf{Decision: } \text{Accept if and only if } v = r$};
        
    \path [line] (sign) -- node[midway, right] {\small Transmit $(M, r, s)$} (ver);
\end{tikzpicture}
\caption{Digital Signature Algorithm (DSA) Signing and Verification Procedures.}
\label{fig:dsa_flowchart}
\end{figure}

\subsection{Domain Parameter and Key Generation}
\begin{enumerate}
    \item \textbf{Domain Parameters:} Publicly shared triple $(p, q, g)$:
    \begin{itemize}
        \item Modulus prime $p$ of bit length $L$ (e.g., 1024 or 2048 bits).
        \item Subgroup prime $q$ of bit length $N$ (e.g., 160 or 256 bits), such that $q \mid (p - 1)$.
        \item Generator $g = h^{(p - 1)/q} \bmod p$ of order $q$.
    \end{itemize}
    \item \textbf{Private Key:} Choose a random integer $x$ in the range:
    \begin{equation}
    1 \le x \le q - 1
    \end{equation}
    \item \textbf{Public Key:} Compute the group element:
    \begin{equation}
    y = g^x \bmod p
    \end{equation}
\end{enumerate}

\subsection{Signature Generation}
To sign a message $M$:
\begin{enumerate}
    \item Compute the cryptographic hash $z = \text{SHA-256}(M)$, truncated to the leftmost $\min(N, 256)$ bits.
    \item Generate a secure, cryptographically random per-message ephemeral nonce $k$ such that:
    \begin{equation}
    1 \le k \le q - 1
    \end{equation}
    \item Compute the signature component $r$:
    \begin{equation}
    r = (g^k \bmod p) \bmod q
    \end{equation}
    If $r = 0$, pick a fresh random $k$ and restart.
    \item Compute $k^{-1} \bmod q$ using the Extended Euclidean Algorithm.
    \item Compute the signature component $s$:
    \begin{equation}
    s = k^{-1} (z + x \cdot r) \bmod q
    \end{equation}
    If $s = 0$, pick a fresh random $k$ and restart.
    \item The digital signature is the pair $(r, s)$.
\end{enumerate}

\subsection{Signature Verification}
Given received message $M$, signature $(r, s)$, and public key $y$:
\begin{enumerate}
    \item Verify that $0 < r < q$ and $0 < s < q$. If either condition fails, reject the signature.
    \item Compute hash integer $z = \text{SHA-256}(M)$ (truncated to $N$ bits).
    \item Compute modular inverse of $s$:
    \begin{equation}
    w = s^{-1} \bmod q
    \end{equation}
    \item Compute intermediate linear combinations:
    \begin{equation}
    u_1 = (z \cdot w) \bmod q, \quad u_2 = (r \cdot w) \bmod q
    \end{equation}
    \item Compute the verification value $v$:
    \begin{equation}
    v = \left( (g^{u_1} \cdot y^{u_2}) \bmod p \right) \bmod q
    \end{equation}
    \item The signature is valid if and only if $v = r$.
\end{enumerate}

\section{Security Assumptions, Properties, and Known Attacks}
\subsection{Security Hardness}
The security of DSA depends on two hard computational problems:
\begin{enumerate}
    \item Discrete logarithm in the full group $\mathbb{Z}_p^*$ (vulnerable to the Number Field Sieve for index calculus).
    \item Discrete logarithm in the subgroup $G_q$ (vulnerable to Pollard's rho algorithm, with complexity $\mathcal{O}(\sqrt{q})$).
\end{enumerate}

\subsection{Catastrophic Failure via Nonce $k$ Leakage or Reuse}
The most infamous vulnerability in DSA stems from improper generation of the ephemeral nonce $k$.
If the same nonce $k$ is used to sign two distinct messages $M_1$ and $M_2$, an attacker observing the resulting signatures $(r, s_1)$ and $(r, s_2)$ can trivially extract the private key $x$:
\begin{align}
s_1 - s_2 &\equiv k^{-1} (z_1 - z_2) \pmod q \implies k \equiv \frac{z_1 - z_2}{s_1 - s_2} \pmod q \\
x &\equiv \frac{s_1 k - z_1}{r} \pmod q
\end{align}
This catastrophic design flaw occurred in the 2010 Sony PlayStation 3 security breach, where Sony's ECDSA/DSA implementation utilized a static nonce, allowing hackers to compute the master private signing key and permanently compromise the platform. \emph{Mitigation:} Generate nonces deterministically according to RFC 6979 \cite{rfc6979} using HMAC-SHA-256.

\section{Advantages, Limitations, and Modern Applications}
\subsection{Advantages}
\begin{itemize}
    \item \textbf{Compact Signatures:} A 1024-bit DSA produces a 320-bit signature $(r, s)$, saving significant network bandwidth compared to 1024-bit RSA.
    \item \textbf{Rapid Signing:} Generation of $r = (g^k \bmod p) \bmod q$ can be precomputed before the message $M$ is known, enabling ultra-fast real-time signing.
\end{itemize}

\subsection{Limitations}
\begin{itemize}
    \item Slower verification than RSA (requires two modular exponentiations).
    \item Highly sensitive to pseudo-random number generator (PRNG) failures.
    \item Cannot perform encryption.
\end{itemize}

\subsection{Modern Applications}
Historically widely used in the Secure Shell (SSH) protocol, GNU Privacy Guard (GPG), and U.S. federal identity infrastructure. Modern deployments have largely transitioned to its elliptic curve analog, ECDSA.

\section{Complete Worked Example Step-by-Step}
We illustrate step-by-step DSA signing and verification using small prime numbers:
\begin{enumerate}
    \item \textbf{Domain Parameters:}
    \begin{itemize}
        \item Choose prime $p = 59$ and prime $q = 29$. Notice $q \mid (p - 1)$ since $59 - 1 = 58 = 2 \times 29$.
        \item Compute generator $g$: Choose $h = 3$, $g = 3^{(59-1)/29} \bmod 59 = 3^2 \bmod 59 = 9$.
    \end{itemize}
    \item \textbf{Key Generation:}
    \begin{itemize}
        \item Choose private key $x = 7 \in [1, 28]$.
        \item Public key $y = g^x \bmod p = 9^7 \bmod 59 = 4782969 \bmod 59 = 16$.
    \end{itemize}
    \item \textbf{Signing Process:}
    \begin{itemize}
        \item Let message hash integer be $z = 15$.
        \item Choose ephemeral nonce $k = 10 \in [1, 28]$.
        \item Compute $r$:
        \begin{equation}
        r = (g^k \bmod p) \bmod q = (9^{10} \bmod 59) \bmod 29 = 41 \bmod 29 = 12
        \end{equation}
        \item Compute $k^{-1} \bmod q$: $10 \times 3 = 30 \equiv 1 \pmod{29}$, so $k^{-1} = 3$.
        \item Compute $s$:
        \begin{equation}
        s = 3 \times (15 + 7 \times 12) \bmod 29 = 3 \times 99 \bmod 29 = 297 \bmod 29 = 7
        \end{equation}
        \item Signature output: $(r, s) = (12, 7)$.
    \end{itemize}
    \item \textbf{Verification Process:}
    \begin{itemize}
        \item Check $0 < 12 < 29$ and $0 < 7 < 29$ (Valid).
        \item Invert $s$: $w = 7^{-1} \bmod 29 = 25$ (since $7 \times 25 = 175 = 6 \times 29 + 1$).
        \item Compute $u_1$:
        \begin{equation}
        u_1 = (z \cdot w) \bmod q = (15 \times 25) \bmod 29 = 375 \bmod 29 = 27
        \end{equation}
        \item Compute $u_2$:
        \begin{equation}
        u_2 = (r \cdot w) \bmod q = (12 \times 25) \bmod 29 = 300 \bmod 29 = 10
        \end{equation}
        \item Compute exponentiations modulo $p = 59$:
        \begin{align*}
        g^{u_1} \bmod p &= 9^{27} \bmod 59 = 51 \\
        y^{u_2} \bmod p &= 16^{10} \bmod 59 = 17
        \end{align*}
        \item Compute $v$:
        \begin{equation}
        \begin{aligned}
        v &= ((51 \times 17) \bmod 59) \bmod 29 \\
          &= (867 \bmod 59) \bmod 29 = 41 \bmod 29 = 12
        \end{aligned}
        \end{equation}
    \end{itemize}
    Since $v = 12$ and $r = 12$, $v = r$ is satisfied, successfully verifying the signature.
\end{enumerate}

\clearpage
\section{Python Implementation from First Principles}
The complete DSA module in \texttt{src/dsa.py} incorporates FIPS 186-4 parameters and first-principles integer math, as presented in Listing \ref{lst:dsa_code}:
\begin{lstlisting}[language=Python, caption={Core DSA Signing and Verification routines.}, label={lst:dsa_code}]
def dsa_sign(message: bytes, x: int, p, q, g):
    if not (0 < x < q):
        raise ValueError("Private key x must satisfy 0 < x < q")
    z = _hash_to_int(message, q.bit_length())
    while True:
        k = secrets.randbelow(q - 1) + 1
        r = pow_mod(g, k, p) % q
        if r == 0: continue
        k_inv = mod_inverse(k, q)
        s = (k_inv * (z + x * r)) % q
        if s == 0: continue
        return r, s

def dsa_verify(message: bytes, signature, y: int, p, q, g):
    r, s = signature
    if not (1 < y < p) or not (0 < r < q and 0 < s < q): return False
    z = _hash_to_int(message, q.bit_length())
    w = mod_inverse(s, q)
    u1 = (z * w) % q
    u2 = (r * w) % q
    v = ((pow_mod(g, u1, p) * pow_mod(y, u2, p)) % p) % q
    return v == r
\end{lstlisting}

\section{Experimental Evaluation and Correctness Verification}
The implementation was tested on the NIST FIPS 186-4 standard domain parameters ($L=1024, N=160$ bits).

\begin{table}[H]
\centering
\caption{DSA Experimental Performance Benchmarks ($L=1024, N=160$).}
\label{tab:dsa_perf}
\begin{tabular}{lc}
\toprule
\textbf{Benchmark Step / Test Case} & \textbf{Observed Experimental Result} \\
\midrule
Key Generation ($x, y$) & 1.30 ms \\
Document Signing Latency ($r, s$) & 2.641 ms \\
Signature Verification Latency ($v \stackrel{?}{=} r$) & 4.727 ms \\
Valid Signature Verification Status & \textbf{VALID (PASSED)} \\
Tampered Document Attack Resistance & \textbf{REJECTED (PASSED)} \\
\bottomrule
\end{tabular}
\end{table}

The verification time ($4.727$ ms) is approximately $1.79 \times$ longer than the signing time ($2.641$ ms) due to the necessity of evaluating two modular exponentiations ($g^{u_1}$ and $y^{u_2}$) during verification.

\needspace{19\baselineskip}
\subsection*{AI Assistance and Prompt Provenance}
\begin{itemize}
    \item \textbf{AI-Assisted Scope:} Theory draft, DLP proofs, worked example, Listing~\ref{lst:dsa_code}, and Table~\ref{tab:dsa_perf} (paired with \texttt{src/dsa.py}).
    \item \textbf{Session Steps:} Baseline in \textbf{Step 1} (\texttt{P01}--\texttt{P05}); layout and citations in \textbf{Steps 3, 5, 6}; arithmetic, CSPRNG, and benchmarks verified in \textbf{Steps 7--9}.
    \item \textbf{Student Verification:} Checked small worked example ($p=59, q=29, v=r=12$), FIPS 186-4 signing/verification, and tampered rejection in \texttt{evidence/\allowbreak terminal\_output.txt}.
\end{itemize}
\begin{lstlisting}[style=promptstyle, caption={Compact Prompt (R-REP-04) for Chapter 4: DSA.}, label={lst:prompt_ch4}]
Following the project specification (Student: TRAN HUU QUYET, SID: 066206014296, assigned DSA via 6 mod 4 = 2), write the DSA chapter in English LaTeX covering all 8 mandatory sections, including a small worked example, pure Python snippets from first principles, and FIPS 186-4 evaluation. Perform an adversarial check of all intermediate modular steps and align Listing 4.1 with src/dsa.py.
\end{lstlisting}
